% =====================================================================
% PLANTILLA: Quiz -- Área de Matemáticas
% María Mercedes Cantillo · Colegio San Alberto Magno
% ---------------------------------------------------------------------
% 1. Copia este archivo con otro nombre (p. ej. quiz-derivadas.tex) en
%    esta misma carpeta (el estilo y el logo se toman de ../estilo/).
% 2. Cambia los datos de la sección "DATOS DEL QUIZ".
% 3. Compila con:  latexmk -pdf quiz-derivadas.tex
%    (o pdflatex dos veces: así aparecen el total de puntos y de páginas).
% =====================================================================
\documentclass[11pt]{article}
% El estilo y el logo viven en la carpeta compartida plantillas/estilo/
\makeatletter\def\input@path{{../estilo/}}\makeatother
\usepackage{mmcantillo}

% =========================== DATOS DEL QUIZ ===========================
\tipodocumento{Quiz}
\asignatura{Cálculo}          % Cálculo, Física, Trigonometría, ...
\titulo{Derivadas: reglas básicas}
\grado{11°}
\periodo{III}
\tiempo{20 min}
% \anio{2026}                 % por defecto: el año actual
% \mostrarsoluciones          % descomenta para imprimir la clave
% ======================================================================

\begin{document}

\encabezado

\begin{instrucciones}
Responde en el espacio indicado. Muestra el procedimiento completo;
las respuestas sin justificación no tienen puntos.
\end{instrucciones}

\begin{preguntas}

\pregunta[2] Calcula la derivada de $f(x) = 3x^4 - 5x^2 + 7x - 1$.
\begin{solucion}[3cm]
$f'(x) = 12x^3 - 10x + 7$.
\end{solucion}

\pregunta[1] Si $g(x) = \sen x \cdot e^{x}$, entonces $g'(x)$ es:
\begin{opciones*}(2)
  \item $\cos x \cdot e^{x}$
  \item $e^{x}(\sen x + \cos x)$
  \item $e^{x}(\cos x - \sen x)$
  \item $-\cos x \cdot e^{x}$
\end{opciones*}
\begin{solucion}
b), por la regla del producto.
\end{solucion}

\pregunta[2] Usa la gráfica de $f(x) = x^3 - 3x$ para estimar los puntos
donde la recta tangente es horizontal. Verifica tu respuesta con $f'(x)=0$.

\begin{center}
\begin{tikzpicture}
\begin{axis}[mmc, xmin=-2.5, xmax=2.5, ymin=-3, ymax=3,
             xtick={-2,-1,...,2}, ytick={-2,0,2}]
  \addplot+[domain=-2.3:2.3] {x^3 - 3*x};
  \addplot[only marks, mark=*, black, mark size=2pt] coordinates {(-1,2) (1,-2)};
\end{axis}
\end{tikzpicture}
\end{center}
\begin{solucion}[3cm]
$f'(x) = 3x^2 - 3 = 0 \iff x = \pm 1$. Puntos: $(-1,\,2)$ y $(1,\,-2)$.
\end{solucion}

\end{preguntas}

\end{document}
